% ============================================================
\chapter{Processus de Poisson}
\label{ch:poisson}
% ============================================================

Combien d'appels arrivent dans un centre t\'el\'ephonique entre 14h et 15h~? Combien de d\'esint\'egrations radioactives se produisent en une seconde~? Combien de clients entrent dans un magasin en une heure~? Dans chacun de ces sc\'enarios, des \'ev\'enements surviennent al\'eatoirement dans le temps, de mani\`ere ind\'ependante et \`a un taux moyen constant. Le processus qui mod\'elise exactement cette situation est le \emph{processus de Poisson}, l'un des objets les plus fondamentaux de la th\'eorie des probabilit\'es. D\'ecouvert par Sim\'eon Denis Poisson en 1837 sous forme de loi de probabilit\'e, il a \'et\'e \'elev\'e au rang de processus stochastique par les travaux de Palm, Feller et Khintchine dans les ann\'ees 1940.

\section{D\'efinition axiomatique}

\begin{definition}[Processus de Poisson]
\label{def:poisson}
Un processus $(N_t)_{t \geq 0}$ \`a valeurs dans $\N$ est un \emph{processus de
Poisson} d'intensit\'e $\lambda > 0$ si~:
\begin{enumerate}
  \item $N_0 = 0$,
  \item $(N_t)$ a des \emph{accroissements ind\'ependants}~: pour tous
    $0 \leq t_1 < t_2 < \cdots < t_n$, les v.a.\ $N_{t_2} - N_{t_1},
    N_{t_3} - N_{t_2}, \ldots, N_{t_n} - N_{t_{n-1}}$ sont ind\'ependantes,
  \item $(N_t)$ a des \emph{accroissements stationnaires}~:
    $N_{t+s} - N_s \sim \text{Poisson}(\lambda t)$ pour tous $s, t \geq 0$,
  \item les trajectoires de $(N_t)$ sont c\`adl\`ag (continues \`a droite, croissantes,
    \`a sauts de taille $1$).
\end{enumerate}
\end{definition}

\begin{remarque}
Les conditions (1)--(3) suffisent \`a caract\'eriser la loi du processus.
La condition (4) porte sur la r\'egularit\'e des trajectoires et peut \^etre obtenue
par construction.
\end{remarque}

\begin{theoreme}[Existence]
Il existe un processus de Poisson d'intensit\'e $\lambda$ pour tout $\lambda > 0$.
\end{theoreme}

\begin{proof}[Id\'ee de la preuve]
On construit $(N_t)$ \`a partir de temps inter-arriv\'ees i.i.d.\ exponentiels~:
soit $(T_n)_{n \geq 1}$ une suite i.i.d.\ $\mathcal{E}(\lambda)$. On pose
$S_n = T_1 + \cdots + T_n$ (le $n$-i\`eme temps d'arriv\'ee) et
$N_t = \max\{n \geq 0 : S_n \leq t\}$. On v\'erifie que $(N_t)$ satisfait les
axiomes de la d\'efinition~\ref{def:poisson}.
\end{proof}

\section{Propri\'et\'es fondamentales}

\begin{proposition}[Moments]
Pour tout $t \geq 0$~:
\[
  \E[N_t] = \lambda t, \quad \Var(N_t) = \lambda t.
\]
\end{proposition}

\begin{proposition}[Temps inter-arriv\'ees]
Les temps entre sauts cons\'ecutifs $T_n = S_n - S_{n-1}$ (avec $S_0 = 0$) sont
i.i.d.\ de loi $\mathcal{E}(\lambda)$.
\end{proposition}

\begin{proposition}[Temps d'arriv\'ee]
Le $n$-i\`eme temps d'arriv\'ee $S_n = T_1 + \cdots + T_n$ suit une loi Gamma~:
$S_n \sim \text{Gamma}(n, \lambda)$, de densit\'e~:
\[
  f_{S_n}(t) = \frac{\lambda^n t^{n-1}}{(n-1)!} e^{-\lambda t}, \quad t > 0.
\]
\end{proposition}

\begin{formules}[title=\textbf{Propri\'et\'es du processus de Poisson}]
\begin{align*}
  N_{t+s} - N_s &\sim \text{Poisson}(\lambda t) \\
  \E[N_t] &= \lambda t, \quad \Var(N_t) = \lambda t \\
  T_n &\sim \mathcal{E}(\lambda) \text{ (i.i.d.)} \\
  S_n &\sim \text{Gamma}(n, \lambda) \\
  \PP(N_t = k) &= \frac{(\lambda t)^k}{k!} e^{-\lambda t}
\end{align*}
\end{formules}

\section{Construction comme CMTC}

\begin{remarque}
Le processus de Poisson est un cas particulier de processus de naissance pure~:
c'est une CMTC sur $\N$ avec $q_{n,n+1} = \lambda$ pour tout $n$ et
$q_{ij} = 0$ pour $j \neq i + 1$. Le g\'en\'erateur est~:
\[
  (Qf)(n) = \lambda [f(n+1) - f(n)].
\]
\end{remarque}

\section{Propri\'et\'e de Markov et absence de m\'emoire}

\begin{theoreme}[Propri\'et\'e de Markov forte]
Le processus de Poisson satisfait la propri\'et\'e de Markov forte~: pour tout
temps d'arr\^et $\tau$ fini p.s., le processus $(N_{\tau+t} - N_{\tau})_{t \geq 0}$
est un processus de Poisson d'intensit\'e $\lambda$, ind\'ependant de $\mathcal{F}_{\tau}$.
\end{theoreme}

\begin{proposition}[Absence de m\'emoire de l'exponentielle]
La loi exponentielle est la seule loi continue \`a la propri\'et\'e d'absence de
m\'emoire~: si $T \sim \mathcal{E}(\lambda)$, alors
\[
  \PP(T > t + s \mid T > t) = \PP(T > s) = e^{-\lambda s}, \quad \forall s, t \geq 0.
\]
\end{proposition}

\section{Op\'erations sur le processus de Poisson}

\subsection{Superposition}

\begin{theoreme}[Superposition]
\label{thm:superposition}
Soient $(N_t^{(1)})$ et $(N_t^{(2)})$ deux processus de Poisson ind\'ependants
d'intensit\'es $\lambda_1$ et $\lambda_2$. Alors $(N_t^{(1)} + N_t^{(2)})$ est un
processus de Poisson d'intensit\'e $\lambda_1 + \lambda_2$.
\end{theoreme}

\subsection{Am\'ecissage (thinning)}

\begin{theoreme}[Am\'ecissage]
\label{thm:thinning}
Soit $(N_t)$ un processus de Poisson d'intensit\'e $\lambda$. Chaque saut est
ind\'ependamment class\'e en type~1 avec probabilit\'e $p$ et type~2 avec
probabilit\'e $1-p$. Alors les processus de comptage $N_t^{(1)}$ et $N_t^{(2)}$
des types~1 et~2 sont des processus de Poisson ind\'ependants d'intensit\'es
$\lambda p$ et $\lambda(1-p)$.
\end{theoreme}

\begin{proof}
Pour tout $t$, conditionnellement \`a $N_t = n$, les $n$ sauts sont class\'es
ind\'ependamment. Donc $N_t^{(1)} \mid (N_t = n) \sim \text{Bin}(n, p)$.
On calcule~:
\[
  \PP(N_t^{(1)} = k) = \sum_{n=k}^{\infty} \binom{n}{k} p^k (1-p)^{n-k}
  \frac{(\lambda t)^n}{n!} e^{-\lambda t} = \frac{(\lambda pt)^k}{k!} e^{-\lambda pt}.
\]
L'ind\'ependance de $N_t^{(1)}$ et $N_t^{(2)}$ se d\'emontre de mani\`ere analogue.
\end{proof}

\subsection{Statistiques d'ordre}

\begin{theoreme}[Statistiques d'ordre]
\label{thm:ordre}
Conditionnellement \`a $N_t = n$, les $n$ temps de saut $(S_1, \ldots, S_n)$
ont la m\^eme loi que les statistiques d'ordre de $n$ v.a.\ i.i.d.\ uniformes
sur $[0, t]$.
\end{theoreme}

\section{Processus de Poisson compos\'e}

\begin{definition}[Processus de Poisson compos\'e]
Soit $(N_t)_{t \geq 0}$ un processus de Poisson d'intensit\'e $\lambda$ et
$(Y_k)_{k \geq 1}$ une suite i.i.d.\ ind\'ependante de $(N_t)$. Le
\emph{processus de Poisson compos\'e} est
\[
  Z_t = \sum_{k=1}^{N_t} Y_k, \quad t \geq 0.
\]
\end{definition}

\begin{proposition}[Moments du processus compos\'e]
\[
  \E[Z_t] = \lambda t \, \E[Y_1], \quad \Var(Z_t) = \lambda t \, \E[Y_1^2].
\]
\end{proposition}

\begin{exemple}
Un mod\`ele de r\'eclamations en assurance~: les r\'eclamations arrivent selon un
processus de Poisson de taux $\lambda$, chaque r\'eclamation ayant un montant
al\'eatoire $Y_k$ (i.i.d.). Le montant total des r\'eclamations jusqu'au temps $t$
est $Z_t = \sum_{k=1}^{N_t} Y_k$.
\end{exemple}

\section{Processus de Poisson non homog\`ene}

\begin{definition}[Processus de Poisson non homog\`ene]
Un processus $(N_t)_{t \geq 0}$ est un \emph{processus de Poisson non homog\`ene}
de \emph{fonction d'intensit\'e} $\lambda(t) \geq 0$ si~:
\begin{enumerate}
  \item $N_0 = 0$,
  \item $(N_t)$ a des accroissements ind\'ependants,
  \item $N_{t+s} - N_s \sim \text{Poisson}\bigl(\int_s^{t+s} \lambda(u) \, du\bigr)$.
\end{enumerate}
\end{definition}

\begin{proposition}[Changement de temps]
Soit $\Lambda(t) = \int_0^t \lambda(u) \, du$ la \emph{fonction d'intensit\'e
cumul\'ee}. Si $(M_s)_{s \geq 0}$ est un processus de Poisson standard
(d'intensit\'e $1$), alors $N_t = M_{\Lambda(t)}$ est un processus de Poisson
non homog\`ene de fonction d'intensit\'e $\lambda(t)$.
\end{proposition}

\section{Processus de Poisson comme martingale}

\begin{proposition}[Martingale compens\'ee]
\label{prop:poisson_martingale}
Si $(N_t)$ est un processus de Poisson d'intensit\'e $\lambda$, alors~:
\begin{enumerate}
  \item $M_t = N_t - \lambda t$ est une martingale (le \emph{processus de Poisson compens\'e}).
  \item $M_t^2 - \lambda t$ est une martingale.
  \item $e^{\theta N_t - \lambda t(e^{\theta} - 1)}$ est une martingale pour tout $\theta \in \R$.
\end{enumerate}
\end{proposition}

\begin{proof}
Pour le point (1), on calcule~:
\begin{align*}
  \E[M_{t+s} - M_s \mid \mathcal{F}_s]
  &= \E[N_{t+s} - N_s - \lambda t \mid \mathcal{F}_s] \\
  &= \E[N_{t+s} - N_s] - \lambda t = \lambda t - \lambda t = 0.
\end{align*}
\end{proof}

\section{Exercices}

\begin{exercice}
Montrer que si $(N_t)$ est un processus de Poisson d'intensit\'e $\lambda$, alors
pour $s < t$ fix\'es, $\E[N_s \mid N_t] = N_t \cdot s/t$.
\end{exercice}

\begin{exercice}
Des appels t\'el\'ephoniques arrivent selon un processus de Poisson de taux
$\lambda = 5$ par heure. Calculer la probabilit\'e d'avoir au moins $3$ appels
en $30$ minutes.
\end{exercice}

\begin{exercice}
Soit $(N_t)$ un processus de Poisson d'intensit\'e $\lambda$ et $T > 0$.
Montrer que $\E[S_1 \mid N_T = n] = T/(n+1)$ o\`u $S_1$ est le premier temps de saut.
\end{exercice}

\begin{exercice}
Soit $Z_t = \sum_{k=1}^{N_t} Y_k$ un processus de Poisson compos\'e avec
$Y_k \sim \mathcal{E}(\mu)$ et $N_t$ Poisson d'intensit\'e $\lambda$.
Calculer $\E[e^{-\alpha Z_t}]$ pour $\alpha > 0$.
\end{exercice}

\begin{exercice}
Montrer le th\'eor\`eme de superposition~\ref{thm:superposition} en utilisant
les fonctions g\'en\'eratrices des probabilit\'es.
\end{exercice}

\begin{exercice}
Des bus arrivent \`a un arr\^et selon un processus de Poisson d'intensit\'e $\lambda$.
Un voyageur arrive \`a un instant d\'eterministe $t_0 > 0$. Montrer que le temps
d'attente du voyageur est exponentiel $\mathcal{E}(\lambda)$, ind\'ependant du
pass\'e.
\end{exercice}
